% !TeX root = ../偏微分方程笔记.tex
% 由主文档通过 \input 载入；请编译主文档。
\setcounter{chapter}{6} % 保留原书章号。
\chapter{一维波动方程}
\label{chap:wave-one}

波动方程把初始扰动沿特征方向传播。有限区间适合模态展开，无界区间则可以得到直接反映传播结构的公式。

\section{无界弦与达朗贝尔公式}
\begin{theorem}[达朗贝尔公式]
设 $g\in C^2(\mathbb R)$、$h\in C^1(\mathbb R)$。$u_{tt}-c^2u_{xx}=0$、$u(x,0)=g(x)$、$u_t(x,0)=h(x)$ 的唯一 $C^2$ 解为
\begin{equation}
 u(x,t)=\frac{g(x-ct)+g(x+ct)}2+
 \frac1{2c}\int_{x-ct}^{x+ct}h(s)\,ds.
 \label{eq:dalembert}
\end{equation}
\end{theorem}
\begin{proof}
特征坐标 $\xi=x-ct$、$\eta=x+ct$ 把方程化为 $v_{\xi\eta}=0$，故 $u=F(x-ct)+G(x+ct)$。初值给出 $F+G=g$ 与 $-cF'+cG'=h$。对第一式求导，与第二式联立，得
\[
 F'=\frac12g'-\frac1{2c}h,\qquad
 G'=\frac12g'+\frac1{2c}h.
\]
任选基点 $x_*$ 积分，得到 $F(s)=g(s)/2-(2c)^{-1}\int_{x_*}^s h+A$、$G(s)=g(s)/2+(2c)^{-1}\int_{x_*}^s h-A$。代入行波表达式，两个常数抵消，两段积分相减，恰得式\eqref{eq:dalembert}。给定的正则性保证公式可微至所需阶数，直接代入恢复三项条件；上述反推又说明任何经典解均只能取此形式。
\end{proof}
零初速度时，初位移分裂为两列振幅各减半、方向相反的行波。非零初速度则通过区间积分贡献位移。

\begin{example}[如何产生单向行波]
设 $g\in C^2(\mathbb R)$，希望初始形状只以速度 $c$ 向右传播，应怎样选择初速度？
\end{example}
\begin{solve}
目标表达式为 $u(x,t)=g(x-ct)$，所以应取 $h(x)=-cg'(x)$。代入达朗贝尔公式，速度项为
\[
 \frac1{2c}\int_{x-ct}^{x+ct}-cg'(s)\,ds
 =\frac12\bigl(g(x-ct)-g(x+ct)\bigr),
\]
与位移项相加后只剩向右的波。若取 $h=cg'$，则只剩向左的波。相同初位移配以不同初速度，可以产生完全不同的运动；这说明波动方程的两个初始条件各有作用。
\end{solve}

\section{依赖域、影响域与有限传播速度}
点 $(x,t)$ 的值仅使用区间 $[x-ct,x+ct]$ 上的初始数据，称其\emph{依赖域（domain of dependence）}。初始点 $x_0$ 只能影响满足 $|x-x_0|\le ct$ 的点，这构成它的\emph{影响域（domain of influence）}。如果 $g,h$ 均在区间 $[a,b]$ 外为零，则 $u(\cdot,t)$ 在 $[a-ct,b+ct]$ 外为零，这就是有限传播。

若两组数据之差为 $\delta g,\delta h$，达朗贝尔公式给出
\[
 |\delta u(x,t)|\le\|\delta g\|_{L^\infty([x-ct,x+ct])}
 +t\|\delta h\|_{L^\infty([x-ct,x+ct])}.
\]
这同时说明局部数据决定性和有限时间内的连续依赖。

\section{半无界弦与边界反射}
半轴固定端 $u(0,t)=0$ 可将 $g,h$ 奇延拓至全轴后使用式\eqref{eq:dalembert}。奇对称使边界值抵消，反射波因此改变符号。若要求延拓为 $g\in C^2(\mathbb R)$、$h\in C^1(\mathbb R)$，半轴数据需满足 $g(0)=g''(0)=h(0)=0$ 等相应端点条件。较少相容性时，公式仍可在适当弱意义下使用，但不能宣称经过所有反射特征后仍是全局 $C^2$ 解。

自由端 $u_x(0,t)=0$ 使用偶延拓，反射不变号；相应光滑延拓需要 $g'(0)=h'(0)=0$。有限区间固定两端时，先作奇延拓再作 $2\ell$ 周期延拓，即可描述反复反射，也可等价地使用正弦级数。

\section{有限弦的正常模态与拨弦实例}
\label{sec:finite-string}
考虑 $0<x<\ell$ 上固定两端的无源弦，初位移为 $g$、初速度为 $h$。令 $u=X T$ 后，空间问题仍为 $-X''=\lambda X$，时间方程却变成 $T''+c^2\lambda T=0$。记
\[
 \omega_k=\frac{ck\pi}{\ell},\qquad
 b_k=\frac2\ell\int_0^\ell g(x)\sin\frac{k\pi x}{\ell}\,dx,
 \qquad d_k=\frac2\ell\int_0^\ell h(x)\sin\frac{k\pi x}{\ell}\,dx.
\]
初值确定每个振子的两个常数，给出
\begin{equation}
 u(x,t)=\sum_{k=1}^\infty
 \left(b_k\cos\omega_kt+\frac{d_k}{\omega_k}\sin\omega_kt\right)
 \sin\frac{k\pi x}{\ell}.
 \label{eq:finite-wave-modes}
\end{equation}
这里 $d_k$ 展开的是初速度，所以在位移公式中必须除以 $\omega_k$。每个空间正弦因子与相应时间振子合称一个\emph{正常模态（normal mode）}，$\omega_k$ 为角频率，其单位为时间的倒数；普通频率为 $\omega_k/(2\pi)$。最低频率称基频，其余整数倍频率描述更细的振型。

公式的可微性需单独检查。一组方便的充分条件是 $\sum k^2|b_k|<\infty$、$\sum k|d_k|<\infty$，此时至二阶的各导数级数一致收敛，得到经典解。更一般地，$g\in H_0^1(0,\ell)$、$h\in L^2(0,\ell)$ 时得到能量解，具体空间含义见第\ref{chap:evolution-weak}章。与热方程不同，余弦和正弦时间因子没有额外衰减，因此波动不会自动消除初值中的尖角。

\begin{example}[在中点拨起后释放]
设弦在中点被拨起高度 $A>0$，从静止释放，初始形状为
\[
 g(x)=\begin{cases}2Ax/\ell,&0\le x\le\ell/2,\\
 2A(\ell-x)/\ell,&\ell/2\le x\le\ell,
 \end{cases}\qquad h(x)=0.
\]
求其模态展开，并解释哪些频率出现。
\end{example}
\begin{solve}
将积分在中点分开，利用 $\int x\sin(\omega x)dx=-x\cos(\omega x)/\omega+\sin(\omega x)/\omega^2$，得到
\[
 b_k=\frac{8A}{k^2\pi^2}\sin\frac{k\pi}{2},\qquad
 u(x,t)=\frac{8A}{\pi^2}\sum_{k=1}^\infty
 \frac{\sin(k\pi/2)}{k^2}\cos\frac{ck\pi t}{\ell}\sin\frac{k\pi x}{\ell}.
\]
偶数模态系数为零，奇数模态系数交替变号。原始形状关于中点对称，而偶数正弦模态关于中点反对称，这给出了积分结果的结构解释。位移级数绝对一致收敛，初位移具有平方可积的分段常数导数，因此得到能量解；初始中点有尖角，不能把它称为从 $t=0$ 起处处二阶光滑的经典解。
\end{solve}

\section{受迫振动与杜阿梅尔原理}
外力在时刻 $s$ 施加的微小冲量相当于该时刻的初速度。把所有时刻的响应相加，得到\emph{杜阿梅尔原理（Duhamel's principle）}：在例如 $f$ 光滑且空间紧支集的充分条件下，零初值的受迫解为
\begin{equation}
 u_f(x,t)=\frac1{2c}\int_0^t\int_{x-c(t-s)}^{x+c(t-s)}f(y,s)\,dy\,ds.
 \label{eq:wave-duhamel}
\end{equation}
积分域是过去光锥的内部。一般初值解等于式\eqref{eq:dalembert}与式\eqref{eq:wave-duhamel}之和，线性使这种分解成立。

\section{例题与习题}
\begin{example}
固定两端的弦初位移、初速度均为零，外力为 $f(x,t)=F\sin(k\pi x/\ell)\cos(\omega t)$，其中 $F\in\mathbb R$、$\omega\ge0$ 为常数，$k$ 为正整数。求受迫解。
\end{example}
\begin{solve}
令 $u=a(t)\sin(k\pi x/\ell)$、$\omega_k=ck\pi/\ell$，得到 $a''+\omega_k^2a=F\cos(\omega t)$、$a(0)=a'(0)=0$。因此
\[
 a(t)=\begin{cases}
 \displaystyle\frac{F}{\omega_k^2-\omega^2}\bigl(\cos(\omega t)-\cos(\omega_kt)\bigr),&\omega\ne\omega_k,\\
 \displaystyle\frac{F}{2\omega_k}t\sin(\omega_kt),&\omega=\omega_k.
 \end{cases}
\]
共振时出现线性增长的振幅；这是无阻尼模型的结论，不能直接用于含阻尼的长期响应。
\end{solve}
\begin{enumerate}
 \item 令 $g=0$、$h$ 为光滑紧支集函数，求某固定位置在足够晚时的位移，并解释一维传播留下的常数位移。
 \item 用奇延拓求固定端半轴问题，检查反射波的符号。
 \item 对常值源项 $f=1$ 和零初值，由式\eqref{eq:wave-duhamel}求解并直接代回方程。
\end{enumerate}

